%% sec_implementation.tex
%% ─────────────────────────────────────────────────────────────────────────────
%% Section 5 — Implementation: HELAS Adaptation and Hardware Optimisations
%% ─────────────────────────────────────────────────────────────────────────────

\section{Implementation}\label{sec:implementation}

\subsection{Data representation: \texttt{std::complex<double>[6]} to
\texttt{aie::vector<cfloat,8>}}

\mgfive generates HELAS code using \texttt{std::complex<double>}
arrays of length~6.  The AIE vector unit natively operates on 8-wide
single-precision complex vectors (\texttt{aie::vector<cfloat,8>}).  The
mapping is:
\begin{itemize}
  \item Lanes~0--5 hold the six physical wavefunction components.
  \item Lanes~6--7 are used as a \emph{cache} for pre-multiplied values
    $V[6] = i V[4]$ and $V[7] = i V[3]$, computed once per wavefunction
    construction via \texttt{stamp\_vec\_i\_cache()}.  Downstream FFV
    functions read these pre-multiplied values directly, eliminating one
    complex multiply per FFV call.
  \item When wavefunctions are passed over the cascade, lanes~6--7 are
    zero-padded (hardware topology does not support 8-wide writes in a single
    beat); the receiving kernel recomputes the cache if needed.
\end{itemize}

\subsection{Complex division elimination via \texttt{cinv()}}

MadGraph propagator denominators of the form $(p^2 - M^2 + iMW)^{-1}$ are
computed using two real divisions (Smith's method applied to complex division).
The AIE implementation replaces this with a custom \texttt{cinv()} that
computes:
\begin{equation}
  \frac{1}{a + ib} = \frac{a}{a^2+b^2} - i\frac{b}{a^2+b^2}
\end{equation}
using a single \texttt{aie::inv()} call (reciprocal of $a^2+b^2$) followed by
two multiplies.  For each Breit-Wigner denominator this reduces the
\texttt{aie::inv()} call count from~2 to~1 ($-50\%$).

\subsection{Vectorised FFV spinor core: \texttt{ffv\_linear\_core()}}

The dominant pattern in \texttt{FFV1\_1/2} is a $4\times 4$ complex linear
system acting on spinor components.  The AIE implementation extracts this as
\texttt{ffv\_linear\_core()}, computing four spinor outputs as 4-wide
\texttt{aie::mac()} chains:
\begin{equation}
  F_i = \sum_{j}\ C_{ij}\, S_j, \quad i\in\{0,1,2,3\},
\end{equation}
where $C_{ij}$ are coupling-dependent coefficients pre-folded from the
\texttt{V3Comb} combination of gauge boson and Dirac-matrix structure.

\paragraph{\texttt{V3Comb} precomputation.}
The gauge boson $\varepsilon^\mu$ enters the FFV vertex through 12 linear
combinations of its four momentum components.  Grouping these into a
\texttt{V3Comb} struct---computed once per FFV call using scalar arithmetic---
eliminates $4\times$ redundant recomputation across the four spinor output
components.

\subsection{Helicity caching (GGTTG\_HEL\_MODE=2)}

The naive implementation recomputes all 5 external wavefunctions for each of
the 32 helicity configurations:
\[
  32 \times [ \texttt{vxxxxx}(w_{g1}) + \texttt{vxxxxx}(w_{g2}) + \ldots ] =
  160\ \text{wavefunction evaluations per PSP.}
\]
Since each particle can take only hel~$=\pm 1$, the implementation instead
pre-computes 10 wavefunctions (2 helicities $\times$ 5 particles) and selects
them via bit-indexed lookup:
\begin{align}
  \textit{bit}_{k} &= (h \gg k) \mathbin{\&} 1, \notag \\
  w_k &= \textit{bit}_{k} \mathbin{?}\ w_k^{+}\ :\ w_k^{-},
  \quad k\in\{0,\ldots,4\}.
\end{align}
This reduces wavefunction evaluations from 160 to 10 per PSP---a \textbf{16$\times$
reduction}---at the cost of 10 vector registers for the wavefunction cache.

\subsection{Pre-normalised colour matrix}

The final colour reduction sums $|M|^2 = \sum_{ij} C_{ij} \operatorname{Re}
(\jamp_i^* \jamp_j)$ using a $6\times 6$ colour matrix.  All division by the
colour normalisation factor is performed at Python/C++ pre-processing time;
the \aie kernel stores only the pre-normalised constants
\texttt{COLOR\_DIAG\_NORM[6]} and \texttt{COLOR\_OFFDIAG\_NORM[15]}, reducing
the inner loop to 21 fused multiply-adds with zero runtime divisions.  Spin
averaging ($\times 1/256$) is a single scalar multiply applied to the total
$|M|^2$ sum.

\subsection{Selective per-kernel HELAS header compilation}

Each kernel \texttt{\#include}s only the HELAS header(s) for the functions it
uses.  This is enforced at the source level:
\begin{itemize}
  \item \texttt{K1}: includes \texttt{processing\_core\_wfgen.h},
    \texttt{processing\_core\_vvv1p0\_1.h} only.  \texttt{VVV1\_0} is
    explicitly absent.
  \item \texttt{K3}: includes \texttt{processing\_core\_ffv.h}
    (for \texttt{FFV1P0\_3}) and \texttt{processing\_core\_vvv1\_0.h} only.
    \texttt{VVV1P0\_1} is explicitly absent.
\end{itemize}
This \emph{compile-time VVV Exclusion Principle} makes any accidental
co-inclusion a compile error.

\subsection{Packet-stream boundary kernels}

\paragraph{K1 (\texttt{kernel\_k1\_wfgen\_pkt}).}
The kernel reads from \texttt{input\_pktstream}: first, a 32-bit packet header
from which the packet ID (bits~[4:0]) identifies the row within the group;
then, 20 $\times$ 32-bit float words encoding the 5-particle PSP
($4$-momenta of $g_1, g_2, t, \bar{t}, g_3$) with TLAST on the last beat.
The 20 words are read as \texttt{int32\_t} and reinterpreted as \texttt{float}
via \texttt{reinterpret\_cast}.  Cascade output is identical to the non-packet
version.

\paragraph{K4 (\texttt{kernel\_k4\_vvvv\_color\_pkt}).}
The kernel calls \texttt{getPacketid(me2\_out, 0)} to retrieve its assigned
packet ID from the \texttt{pktmerge} infrastructure, writes a packet header
via \texttt{writeHeader(me2\_out, pktType, ID)}, then writes the
\texttt{BATCH} ME$^2$ results as 32-bit floats with TLAST on the final beat.
Cascade input processing is identical to the non-packet version.
